\chapter{Implementación}
\label{chap:implem}


\section{Implementación da API}
%meter en apéndice
% - Exemplos do JSON devolto por cada endpoint,, exemplos de código e igual unha foto do swagger e de algunha chamada.
% - Os modelos de pydantic
% - Mteer igual cousas mais concretas dos edpoint auxiliares


A \gls{api} implementouse con FastAPI e mantén separadas a configuración dos servizos externos, os puntos de entrada, os esquemas de datos e a lóxica auxiliar. 

\subsection{Inicialización e acceso aos servizos externos}

A conexión cos servizos externos realízase durante o arranque da \gls{api}, antes de que esta comece a atender peticións. Para iso emprégase o mecanismo de ciclo de vida de \textit{FastAPI}, no que se inicializan \textit{Azure Document Intelligence} e \textit{Firebase Admin} e se recuperan as credenciais necesarias para acceder a estes servizos.

A obtención das credenciais segue unha estratexia común nos dous casos. En primeiro lugar, compróbase se están dispoñibles nas variables de contorno do sistema. Este mecanismo facilita a execución da \gls{api} durante o desenvolvemento sen necesidade de acceder a servizos adicionais. Se as credenciais non están presentes, utilízase \textit{DefaultAzureCredential} para autenticarse contra \textit{Azure Key Vault} e recuperar os segredos correspondentes.

Unha vez obtida a clave de \textit{Azure Document Intelligence}, créase unha instancia de \textit{DocumentIntelligenceClient}. Esta instancia almacénase no estado da aplicación e reutilízase durante todo o seu ciclo de vida. As rutas que necesitan realizar a extracción reciben o cliente mediante o sistema de inxección de dependencias de \textit{FastAPI}, evitando así ter que crealo e configuralo para cada petición.

No caso de \textit{Firebase Admin}, as credenciais recuperadas conteñen a configuración da conta de servizo. A partir delas inicialízase o SDK de administración de Firebase. Antes de realizar esta operación compróbase se xa existe unha instancia inicializada, evitando rexistrar a aplicación máis dunha vez durante a execución. Unha vez configurado, o servizo de mensaxería de Firebase pode empregarse desde as rutas encargadas do envío de notificacións.

%Esta implementación centraliza a configuración dos servizos externos no arranque da \gls{api}, mantén a xestión das credenciais separada da lóxica das rutas e permite reutilizar os recursos xa inicializados. Ademais, a combinación de variables de contorno e \textit{Azure Key Vault} facilita o uso de configuracións diferentes segundo o contorno de execución sen incorporar credenciais directamente no código fonte.

Esta implementación centraliza a configuración dos servizos externos e mantén a xestión das credenciais separada da lóxica dos endpoints.

\subsection{Implementación do procesamento dunha folla}
A solución seleccionada no Capítulo~\ref{chap:procesamento-follas} integrouse no endpoint \textit{/extraccion/}. Ademais de  invocar ao servizo de procesamento documental, o endpoint executa as operacións necesarias para preparar e comprobar os resultados. %Meter preprocesamento

%TODO-> Meter no anexo como funciona ben o preprocesamento (Como se implementa). Revisar capitulo de procesamento 
Antes de realizar a extracción compróbase que o ficheiro se corresponda con algún dos formatos admitidos. A continuación aplícase a etapa de preprocesamento e envíase o resultado ao modelo personalizado de Azure. Se o servizo non identifica ningún documento ou a confianza global da extracción é inferior ao 60\%, a petición considérase non válida e non se constrúe a pauta.

Unha vez obtidos os campos, realízanse distintas comprobacións para evitar resultados incoherentes. A data da próxima visita compárase coa data do informe e rexéitase a extracción se é anterior. O valor do INR normalízase para admitir tanto coma como punto decimal, compróbase que poida converterse a un valor numérico, que se atope dentro dun intervalo lóxico de 0,5 a 10 e que a confianza proporcionada por Azure sexa polo menos do 80\,\%.

A táboa de doses procésase cela a cela para obter a data, a dose e a acción correspondente. Durante esta transformación compróbase que o día e o mes permitan construír unha data válida e tense en conta o cambio de ano cando a pauta pasa de decembro a xaneiro. As celas que indican \textit{NO TOMAR} convértense nunha dose cero, mentres que as identificadas como \textit{CONTROL} reciben un tratamento específico. Se non é posible recuperar correctamente a data da cela de control, emprégase a data da próxima visita. Do mesmo xeito, se a data extraída non coincide coa próxima visita, substitúese por esta para manter a coherencia entre ambos campos. Finalmente, se non se consegue recuperar ningunha entrada válida da táboa, a extracción é rexeitada.

O histórico transfórmase ao formato interno conservando unicamente as filas nas que se recuperou polo menos a data ou o INR. Unha vez procesada e validada a información, os resultados organízanse empregando os modelos de Pydantic cos que se implementou a estrutura definida no deseño. A resposta constrúese mediante \textit{AnalisisResponse}, garantindo que os datos enviados ao cliente móbil manteñan o formato establecido pola \gls{api}.

\subsection{Endpoints auxiliares}
Ademais do procesamento das follas, a \gls{api} incorpora varios endpoints auxiliares. O endpoint de centros sanitarios permite buscar os datos dun centro a partir do seu nome, empregando o catálogo almacenado na \gls{api} e admitindo pequenas variacións na denominación. O endpoint de estado permite comprobar que a \gls{api} está operativa e devolve o estado, a data e hora da consulta e a versión en execución, mentres que o de preprocesamento permite aplicar esta etapa de forma independente, sen realizar a extracción con Azure, facilitando así a comprobación do resultado xerado.

A \gls{api} tamén intervén na comunicación entre os dispositivos vinculados mediante o endpoint \textit{/notificar/enviar}. Este recibe o token do dispositivo destinatario, o tipo de aviso e un \textit{payload} que xa foi cifrado no dispositivo de orixe. A API non necesita interpretar o seu contido, senón que actúa como intermediaria e emprega Firebase Cloud Messaging para reenvialo ao destinatario. Para determinados avisos engádese unha notificación visible cun título e unha mensaxe xenéricos, mantendo a información sensible exclusivamente no \textit{payload} cifrado.

\subsection{Tratamento de erros}
Diferéncianse os erros segundo a súa orixe. Os ficheiros ou resultados que non superan as comprobacións da extracción devolven un erro 400 \textit{(Bad Request)}, mentres que os erros de validación dos parámetros de entrada son xestionados por FastAPI mediante unha resposta 422\textit{ (Unprocessable Entity)}. Os fallos de comunicación con Azure devolven un 502 \textit{(Bad Gateway)} e os erros inesperados da aplicación un 500 \textit{(Internal Server Error)}. Noutros endpoints empréganse tamén códigos específicos, como o 404\textit{ (Not Found)} cando non se localiza un centro sanitario ou o 503 \textit{(Service Unavailable)} cando non se pode acceder ao catálogo. Esta diferenciación permite que o cliente identifique mellor a causa do problema e mostre unha mensaxe adecuada á persoa usuaria.

\section{Implementación do cliente móbil}

%Revisa rintroducció
%A aplicación móbil implementouse en Flutter seguindo o patrón MVVM definido no deseño. As vistas xestionan a interacción coa persoa usuaria, os modelos de vista manteñen o estado e coordinan as operacións, e os modelos representan os datos empregados pola aplicación. Os servizos encapsulan o acceso á \gls{api}, á persistencia, ás notificacións e á sincronización.

\subsection{Inicialización da aplicación}

Antes de mostrar a interface, a aplicación inicializa os servizos necesarios para o seu funcionamento mediante a clase \textit{InicializadorAplicacion}.

Neste proceso configúrase Firebase, prepárase a recepción de mensaxes tanto en primeiro como en segundo plano e solicítanse os permisos necesarios para as notificacións. Tamén se recupera do almacenamento seguro a información da configuración previa, como o rol seleccionado e se o proceso inicial xa foi completado.

A partir destes datos determínase a pantalla que debe mostrarse ao iniciar a aplicación. Se a configuración está pendente, accédese ao fluxo inicial. Se xa foi completada, móstrase directamente a interface de paciente ou de persoa supervisora. No caso do paciente, tamén se restauran os recordatorios asociados á pauta activa.

\subsection{Incorporación e revisión dunha folla de tratamento}
A incorporación dunha folla de tratamento divídese en 3 pasos: o envñio á \gls{api}, a revisión dos resultados e a confirmación final. O ficheiro pode obterse mediante a cámara, a galería ou o selector de documentos. O modelo de vista de captura mantén o estado do proceso e xestiona os posibles erros.

O envío realízase mediante \textit{ApiService}, empregando unha petición \gls{http} de tipo \textit{multipart}. Se a extracción é correcta, o \gls{json} devolto transfórmase nun \textit{AnaliseModel}. Se se produce un erro, a mensaxe recibida trasládase á interface para informar á persoa usuaria.

Antes de gardar a pauta, os datos pasan por unha fase de revisión xestionada por \textit{RevisionPauta}. Esta permite modificar, engadir ou eliminar días, cambiar doses e marcar días de control. Tamén comproba que a pauta non estea baleira, que as datas sexan válidas e non estean duplicadas, e que as doses teñan un formato admitido.

Tras a confirmación, o comportamento depende do dispositivo. Se a folla foi incorporada polo paciente, gárdase como pauta activa, reprograman os recordatorios e notifícase o cambio ás persoas supervisoras vinculadas. Se foi incorporada desde o dispositivo supervisor, a pauta revisada envíase ao paciente para que sexa almacenada como tratamento activo.





